%% ============================================================================
%% Section: Related Work
%% Included from main.tex via %% ============================================================================
%% Section: Related Work
%% Included from main.tex via %% ============================================================================
%% Section: Related Work
%% Included from main.tex via %% ============================================================================
%% Section: Related Work
%% Included from main.tex via \input{sections/related}
%% Citations use placeholder keys in refs.bib (TODO_*) -- verify before CR.
%% ============================================================================
\section{Related Work}\label{sec:related}

\paragraph{Reachability indexing}
Answering reachability without per-query traversal is a long line of database
work, from 2-hop labeling~\cite{TODO_2hop} through scalable indexes such as
GRAIL~\cite{TODO_grail}, pruned landmark labeling~\cite{TODO_pll}, and
topological-folding labels~\cite{TODO_tflabel}. Our ancestor/descendant bitsets
are a deliberately simple point in this space---a full transitive closure packed
into machine words---chosen because causal DAGs are comparatively small and the
queries above it (d-separation, adjustment enumeration) need whole ancestor and
descendant \emph{sets}, not just point-to-point reachability, which bitsets
deliver as one word-parallel operation. Positioning Causal4j in this lineage is
the point: we bring a classical reachability-indexing trade-off to a workload
(causal structural queries) that has not previously been treated as an indexed
DBMS workload.

\paragraph{Property-graph query languages}
Cypher and the emerging GQL/SQL-PGQ standards~\cite{TODO_gql} express graph
pattern matching and variable-length paths declaratively, and our language
compiles to Cypher plus procedures rather than replacing it. What Cypher does not
provide is conditioning-aware d-separation or adjustment-set enumeration as
first-class operators; these are not path patterns, and expressing them in Cypher
means client-side orchestration of many traversals. Causal4j adds the causal
operators while keeping the model queryable by ordinary Cypher.

\paragraph{Causal tooling}
Causal analysis is dominated by standalone libraries---DoWhy~\cite{TODO_dowhy},
dagitty~\cite{TODO_dagitty}, pgmpy~\cite{TODO_pgmpy},
causal-learn~\cite{TODO_causallearn}, and Ananke/Y0~\cite{TODO_ananke}---that
implement d-separation, the backdoor criterion, and identification over an
in-memory graph. None is DBMS-resident: each requires exporting the DAG, computes
in a separate process, and discards its representation, so none offers durable
shared access, in-database querying, or incremental maintenance under a changing
graph. We reuse their semantics (indeed we validate against them in
\S\ref{sec:eval}) but relocate the computation inside the store.

\paragraph{In-DBMS analytics and incremental maintenance}
Pushing analytics into the database via procedures and user-defined functions is
established practice; our contribution is the specific physical design and query
processing for causal operators, not the extension mechanism. Finally, keeping the
index consistent under updates connects to incremental view and index
maintenance~\cite{TODO_ivm}: our affected-cone patch of the reachability closures
(\S\ref{sec:incr}) is an instance of maintaining a materialized transitive
closure under edge insertion, with a principled fallback to recomputation for the
non-local case of deletion.

%% Citations use placeholder keys in refs.bib (TODO_*) -- verify before CR.
%% ============================================================================
\section{Related Work}\label{sec:related}

\paragraph{Reachability indexing}
Answering reachability without per-query traversal is a long line of database
work, from 2-hop labeling~\cite{TODO_2hop} through scalable indexes such as
GRAIL~\cite{TODO_grail}, pruned landmark labeling~\cite{TODO_pll}, and
topological-folding labels~\cite{TODO_tflabel}. Our ancestor/descendant bitsets
are a deliberately simple point in this space---a full transitive closure packed
into machine words---chosen because causal DAGs are comparatively small and the
queries above it (d-separation, adjustment enumeration) need whole ancestor and
descendant \emph{sets}, not just point-to-point reachability, which bitsets
deliver as one word-parallel operation. Positioning Causal4j in this lineage is
the point: we bring a classical reachability-indexing trade-off to a workload
(causal structural queries) that has not previously been treated as an indexed
DBMS workload.

\paragraph{Property-graph query languages}
Cypher and the emerging GQL/SQL-PGQ standards~\cite{TODO_gql} express graph
pattern matching and variable-length paths declaratively, and our language
compiles to Cypher plus procedures rather than replacing it. What Cypher does not
provide is conditioning-aware d-separation or adjustment-set enumeration as
first-class operators; these are not path patterns, and expressing them in Cypher
means client-side orchestration of many traversals. Causal4j adds the causal
operators while keeping the model queryable by ordinary Cypher.

\paragraph{Causal tooling}
Causal analysis is dominated by standalone libraries---DoWhy~\cite{TODO_dowhy},
dagitty~\cite{TODO_dagitty}, pgmpy~\cite{TODO_pgmpy},
causal-learn~\cite{TODO_causallearn}, and Ananke/Y0~\cite{TODO_ananke}---that
implement d-separation, the backdoor criterion, and identification over an
in-memory graph. None is DBMS-resident: each requires exporting the DAG, computes
in a separate process, and discards its representation, so none offers durable
shared access, in-database querying, or incremental maintenance under a changing
graph. We reuse their semantics (indeed we validate against them in
\S\ref{sec:eval}) but relocate the computation inside the store.

\paragraph{In-DBMS analytics and incremental maintenance}
Pushing analytics into the database via procedures and user-defined functions is
established practice; our contribution is the specific physical design and query
processing for causal operators, not the extension mechanism. Finally, keeping the
index consistent under updates connects to incremental view and index
maintenance~\cite{TODO_ivm}: our affected-cone patch of the reachability closures
(\S\ref{sec:incr}) is an instance of maintaining a materialized transitive
closure under edge insertion, with a principled fallback to recomputation for the
non-local case of deletion.

%% Citations use placeholder keys in refs.bib (TODO_*) -- verify before CR.
%% ============================================================================
\section{Related Work}\label{sec:related}

\paragraph{Reachability indexing}
Answering reachability without per-query traversal is a long line of database
work, from 2-hop labeling~\cite{TODO_2hop} through scalable indexes such as
GRAIL~\cite{TODO_grail}, pruned landmark labeling~\cite{TODO_pll}, and
topological-folding labels~\cite{TODO_tflabel}. Our ancestor/descendant bitsets
are a deliberately simple point in this space---a full transitive closure packed
into machine words---chosen because causal DAGs are comparatively small and the
queries above it (d-separation, adjustment enumeration) need whole ancestor and
descendant \emph{sets}, not just point-to-point reachability, which bitsets
deliver as one word-parallel operation. Positioning Causal4j in this lineage is
the point: we bring a classical reachability-indexing trade-off to a workload
(causal structural queries) that has not previously been treated as an indexed
DBMS workload.

\paragraph{Property-graph query languages}
Cypher and the emerging GQL/SQL-PGQ standards~\cite{TODO_gql} express graph
pattern matching and variable-length paths declaratively, and our language
compiles to Cypher plus procedures rather than replacing it. What Cypher does not
provide is conditioning-aware d-separation or adjustment-set enumeration as
first-class operators; these are not path patterns, and expressing them in Cypher
means client-side orchestration of many traversals. Causal4j adds the causal
operators while keeping the model queryable by ordinary Cypher.

\paragraph{Causal tooling}
Causal analysis is dominated by standalone libraries---DoWhy~\cite{TODO_dowhy},
dagitty~\cite{TODO_dagitty}, pgmpy~\cite{TODO_pgmpy},
causal-learn~\cite{TODO_causallearn}, and Ananke/Y0~\cite{TODO_ananke}---that
implement d-separation, the backdoor criterion, and identification over an
in-memory graph. None is DBMS-resident: each requires exporting the DAG, computes
in a separate process, and discards its representation, so none offers durable
shared access, in-database querying, or incremental maintenance under a changing
graph. We reuse their semantics (indeed we validate against them in
\S\ref{sec:eval}) but relocate the computation inside the store.

\paragraph{In-DBMS analytics and incremental maintenance}
Pushing analytics into the database via procedures and user-defined functions is
established practice; our contribution is the specific physical design and query
processing for causal operators, not the extension mechanism. Finally, keeping the
index consistent under updates connects to incremental view and index
maintenance~\cite{TODO_ivm}: our affected-cone patch of the reachability closures
(\S\ref{sec:incr}) is an instance of maintaining a materialized transitive
closure under edge insertion, with a principled fallback to recomputation for the
non-local case of deletion.

%% Citations use placeholder keys in refs.bib (TODO_*) -- verify before CR.
%% ============================================================================
\section{Related Work}\label{sec:related}

\paragraph{Reachability indexing}
Answering reachability without per-query traversal is a long line of database
work, from 2-hop labeling~\cite{TODO_2hop} through scalable indexes such as
GRAIL~\cite{TODO_grail}, pruned landmark labeling~\cite{TODO_pll}, and
topological-folding labels~\cite{TODO_tflabel}. Our ancestor/descendant bitsets
are a deliberately simple point in this space---a full transitive closure packed
into machine words---chosen because causal DAGs are comparatively small and the
queries above it (d-separation, adjustment enumeration) need whole ancestor and
descendant \emph{sets}, not just point-to-point reachability, which bitsets
deliver as one word-parallel operation. Positioning Causal4j in this lineage is
the point: we bring a classical reachability-indexing trade-off to a workload
(causal structural queries) that has not previously been treated as an indexed
DBMS workload.

\paragraph{Property-graph query languages}
Cypher and the emerging GQL/SQL-PGQ standards~\cite{TODO_gql} express graph
pattern matching and variable-length paths declaratively, and our language
compiles to Cypher plus procedures rather than replacing it. What Cypher does not
provide is conditioning-aware d-separation or adjustment-set enumeration as
first-class operators; these are not path patterns, and expressing them in Cypher
means client-side orchestration of many traversals. Causal4j adds the causal
operators while keeping the model queryable by ordinary Cypher.

\paragraph{Causal tooling}
Causal analysis is dominated by standalone libraries---DoWhy~\cite{TODO_dowhy},
dagitty~\cite{TODO_dagitty}, pgmpy~\cite{TODO_pgmpy},
causal-learn~\cite{TODO_causallearn}, and Ananke/Y0~\cite{TODO_ananke}---that
implement d-separation, the backdoor criterion, and identification over an
in-memory graph. None is DBMS-resident: each requires exporting the DAG, computes
in a separate process, and discards its representation, so none offers durable
shared access, in-database querying, or incremental maintenance under a changing
graph. We reuse their semantics (indeed we validate against them in
\S\ref{sec:eval}) but relocate the computation inside the store.

\paragraph{In-DBMS analytics and incremental maintenance}
Pushing analytics into the database via procedures and user-defined functions is
established practice; our contribution is the specific physical design and query
processing for causal operators, not the extension mechanism. Finally, keeping the
index consistent under updates connects to incremental view and index
maintenance~\cite{TODO_ivm}: our affected-cone patch of the reachability closures
(\S\ref{sec:incr}) is an instance of maintaining a materialized transitive
closure under edge insertion, with a principled fallback to recomputation for the
non-local case of deletion.
